\paragraph{\trLineBundle{}上的 \texorpdfstring{$\bar{\partial}$}{dbar}-算子。}%
\label{sub:n-operators-line-bundles}
在闭\trRiemann\trSurface{}上，负\trDegree{}的\trLineBundle{}不可能有非平凡的%
\trHolomorphic\trSection。有一个标准论证能把这一事实推广到带边的\trRiemann\trSurface，以及
形如 $\bar\partial+\{\text{零阶项}\}$ 的算子。具体地说，设 $S$ 是一个带非空边界的%
\trCompact\trConnected\trRiemann\trSurface；$E\to S$ 是一个\trHermitian\trLineBundle；$F\subset
E|\partial S$ 是一个实\trSubbundle；而 $\bar\partial_\nabla$ 是与之相关联的一个
$\bar\partial$-算子。在选取 $E$ 的一个\trTrivialization{}之后，可以把 $F$ 看成有限多个
映射 $\rho_1,\ldots,\rho_l\colon S^1\to\RP^1$，每个\trBoundaryComponent{}对应一个。以
$\mu(\rho_j)$ 记它们的\trDegree。假设 $\bar\partial_\nabla X=0$ 有一个非零解。它的诸零点
是\trIsolated{}的，每个都有有限阶 $\nu(z)$，而一个简单的\textbf{卷绕数}（winding number）
论证表明
\begin{equation}\label{eq:winding-number-count}
2\sum_{z\in\mathrm{int}(S)}\nu(z)+\sum_{z\in\partial S}\nu(z)
=\mu(\rho_1)+\cdots+\mu(\rho_l).
\end{equation}
特别地，
\begin{lemma}\label{thm:injective-negative-degree}
若 $\mu(\rho_1)+\cdots+\mu(\rho_l)<0$，则 $\bar\partial_\nabla$ 是\trInjective。
\noqed
\end{lemma}
对带\trStriplikeEnd{}的\trSurface{} $S$ 及其上的\trNondegenerate{} $\bar\partial_\nabla$-算子，
本结果还有一个版本。为简单起见，假设在每个端 $\epsilon_\zeta\colon Z^\pm\to S$ 上，
$(E,\nabla)$ 的\trPullback{}是平凡的，并且边界\trSubbundle{}是\trLocallyConstant{}的。必要时把%
\trTrivialization{} $\epsilon_\zeta^*E\cong Z^\pm\times\CC$ 旋转一个常数，便可写成
\begin{equation}\label{eq:end-boundary-lines}
F_{\epsilon_\zeta(s,0)}=\RR,\qquad
F_{\epsilon_\zeta(s,1)}=e^{i\sigma_\zeta}\RR
\end{equation}
对某个 $\sigma_\zeta\in(-\pi;0)$。现在仍可像前面那样得到诸映射 $\rho_j$，即对%
\trCompactification{} $\hat{S}$ 的每个\trBoundaryComponent{}取一个：沿 $F$ 追踪，并且在每个
无穷远点附近，取\eqref{eq:end-boundary-lines} 中两条直线之间那条经过
$e^{it\sigma_\zeta}\RR$（$t\in[0;1]$）的\trHomotopy。
\begin{lemma}\label{thm:winding-injectivity}
若 $\mu(\rho_1)+\cdots+\mu(\rho_l)-|\Sigma^-|<0$，则 $\bar\partial_\nabla$ 是\trInjective。
\end{lemma}
为证此，考虑 $\bar\partial_\nabla X=0$ 的一个非零解。在每个\trStriplikeEnd{}上，该算子
就是带\trBoundaryCondition{} \eqref{eq:end-boundary-lines} 的标准
$\bar\partial_s+i\partial_t$，于是 $X_\zeta(s,t)=X(\epsilon_\zeta(s,t))$ 可以作%
\trFourier{}展开：
\begin{equation*}
X_\zeta(s,t)=
\begin{cases}
\sum_{m>0}c_{\zeta,m}e^{(\sigma_\zeta+\pi m)z}, & \zeta\in\Sigma^-,\\
\sum_{m\ge 0}c_{\zeta,m}e^{(\sigma_\zeta-\pi m)z}, & \zeta\in\Sigma^+,
\end{cases}
\end{equation*}
其中\trCoefficients{} $c_{\zeta,m}\in\RR$ 在 $m\to\infty$ 时快速衰减。以 $\nu(\zeta)$ 记使
$c_{\zeta,m}\neq 0$ 的最小的 $m$。对充分大的 $|s|$，这确定了非零函数
$X_\zeta|_{\{s\}\times[0;1]}$ 的\trTopology{}性态，而对\eqref{eq:winding-number-count} 所
产生的相应修正项给出
\begin{equation*}
2\sum_{z\in\mathrm{int}(S)}\nu(z)+\sum_{z\in\partial S}\nu(z)+\sum_{\zeta\in\Sigma}\nu(\zeta)
=\mu(\rho_1)+\cdots+\mu(\rho_l).
\end{equation*}
由于对 $\zeta\in\Sigma^-$ 有 $\nu(\zeta)>0$，所欲证的结果得证。
\begin{remark}\label{thm:zeros-finite-multiplicity}
$\bar\partial_\nabla X=0$ 的非平凡解只有\trIsolated{}零点、且每个零点都有有限%
\trMultiplicity，这一事实是 Aronszajn\footnote{译注：原书此处人名印作 Aronszaijn；疑应为 Aronszajn（人名标准拼法）；此处径改。——译者注}\trUniqueContinuation{}定理的一个推论：%
\trInteriorPoint{}情形见 [49]，边界点情形见 [91]。同一定理（或更直接的\trTaylor{}级数计算）还表明局部\trMultiplicity{}都是
正的。最后应当指出，即使去掉在每个\trStriplikeEnd{}上平凡性的假设，
\cref{thm:winding-injectivity} 仍然成立。证明的原理不变，但 $X_\zeta$ 不再能被显式地
写出，于是 $\nu(\zeta)$ 的定义要转而依赖对 $X_\zeta$ 的\trAsymptotic{}性态的更精致的
分析；见 [114]。
\end{remark}
